\subsubsection{HTML-এর সীমাবদ্ধতা: কোন কাজের জন্য বাড়তি প্রযুক্তি লাগে?}
HTML বিষয়বস্তুর পরিচয় ও গঠন দেয়। শুধু মার্কআপ লিখে সব ধরনের কাজ করা যায় না।
\begin{enumerate}
\item \textbf{সূক্ষ্ম নকশা:} রং ও বিন্যাস নিয়ন্ত্রণে CSS প্রয়োজন।
\item \textbf{নিজস্ব হিসাব ও প্রতিক্রিয়া:} নিজের নিয়মে তাৎক্ষণিক হিসাব বা প্রতিক্রিয়ার জন্য JavaScript প্রয়োজন।
\item \textbf{সার্ভারের কাজ:} জমা পড়া আবেদন যাচাই, সংরক্ষণ ও পরে ফেরত দিতে সার্ভারের প্রোগ্রাম প্রয়োজন।
\item \textbf{প্রোগ্রাম চালানোর সীমা:} HTML-এর নিজস্ব শর্ত, লুপ বা ডেটাবেজে তথ্য সংরক্ষণের প্রোগ্রাম নেই।
\item \textbf{হাতে সম্পাদনার শ্রম:} বড় নথি হাতে সম্পাদনা শ্রমসাধ্য; প্রয়োজন অনুযায়ী অন্য প্রযুক্তির সঙ্গে HTML ব্যবহার করা হয়।
\end{enumerate}
\textbf{HTML-এর কাজ:} ভর্তি-ফরমের ঘর, বোতাম, নির্দেশনা ও ফি-তালিকা দেখানো যায়। \texttt{required}-এর মতো অ্যাট্রিবিউট দিয়ে ব্রাউজারের কিছু সাধারণ যাচাইও চাওয়া যায়।
\textbf{মনে রাখবে:} HTML ব্যবহৃত হয়েছে বলেই কোনো সাইট সব সময় স্থির বা অনিরাপদ হবে—এমন সিদ্ধান্ত নেওয়া যাবে না। HTML-এ ফি-তালিকা দেখানো, আবেদনকারীর ফি গণনা করা ও ফল সংরক্ষণ করা আলাদা কাজ। তালিকা দেখালেই হিসাব ও সংরক্ষণ স্বয়ংক্রিয় হয় না।
